\subsection{特征函数展开}\label{sec:elliptic-eigenfunction-expansions}

在椭圆算子 $L$ 的适当假设下, Theorem~\ref{thm:weak-dirichlet-positive-shift-existence} 保证存在 $\lambda_G$, 使得若 $\widetilde\lambda>\lambda_G$, 则对任意 $f\in H^{-k}(\Omega)$, 下列方程的 Dirichlet 问题存在唯一的 (弱) 解 $u\in H_0^k(\Omega)$:\[\widetilde L(\bm{x},D)u:=L(\bm{x},D)u+\widetilde\lambda u=f.\]本节将把前一章中发展的一些算子方法用于这一问题. 这将使现有的存在性理论在两个基本方面得到改进. 首先, Fredholm 定理将提供 $\widetilde\lambda<\lambda_G$ 时解的存在性与唯一性的信息. 其次, 若算子 $L$ 满足某种对称性条件, 则可用特征函数展开法构造解 (在实际中则是近似解).

\subsubsection{Fredholm 理论}\label{subsec:elliptic-fredholm-theory}

本节考虑 $f\in L^2(\Omega)$ 时的非齐次特征值问题
\begin{equation}\label{eq:elliptic-nonhomogeneous-shifted-eigenvalue-problem}
L(\bm{x},D)u+\lambda u=f
\end{equation}
其中 $L(\bm{x},D)$ 为算子\[L(\bm{x},D)u=\sum_{0\leq|\sigma|,|\gamma|\leq k}(-1)^{|\sigma|}D^\sigma(a_{\sigma\gamma}(\bm{x})D^\gamma u),\]与 $L$ 相联系的双线性型为\[B[v,u]=\sum_{0\leq|\sigma|,|\gamma|\leq k}\int_\Omega a_{\sigma\gamma}(\bm{x})D^\gamma uD^\sigma v\,\mathrm{d}\bm{x}.\]

假设 Theorem~\ref{thm:weak-dirichlet-positive-shift-existence} 的条件成立, 固定 $\widetilde\lambda>\lambda_G$, 且 $\widetilde\lambda>0$. 那么对任意 $f\in L^2(\Omega)$, 下列问题存在唯一解 $u\in H_0^k(\Omega)$:
\begin{equation}\label{eq:elliptic-shifted-solution-variational-definition}
B_{\widetilde\lambda}[v,u]:=B[v,u]+\widetilde\lambda(v,u)=(v,f)_{L^2(\Omega)}\quad\text{for every }v\in H_0^k(\Omega).
\end{equation}
现在如下定义算子 $\overline G:L^2(\Omega)\to H_0^k(\Omega)$: 对每个 $f\in L^2(\Omega)$, 定义
\begin{equation}\label{eq:elliptic-sobolev-solution-operator-definition}
\overline G(f):=\widetilde\lambda u,
\end{equation}
其中 $u$ 为下列方程的 Dirichlet 问题的唯一 (弱) 解:
\begin{equation}\label{eq:elliptic-solution-operator-shifted-dirichlet-equation}
L(\bm{x},D)u+\widetilde\lambda u=f;
\end{equation}
即 $u$ 解 \eqref{eq:elliptic-shifted-solution-variational-definition}. 换言之, 对每个 $f\in L^2(\Omega)$ 和 $v\in H_0^k(\Omega)$, 有
\begin{equation}\label{eq:elliptic-sobolev-solution-operator-variational-identity}
B_{\widetilde\lambda}[v,\overline G(f)]=\widetilde\lambda(v,f)_{L^2(\Omega)}.
\end{equation}


形式上有
\begin{equation}\label{eq:elliptic-sobolev-solution-operator-formal-inverse}
\overline G=\widetilde\lambda(L+\widetilde\lambda)^{-1}.
\end{equation}
由 \eqref{eq:weak-dirichlet-solution-negative-norm-bound}, 算子 $\overline G$ 有界. 现在以 $\overline G$ 和 $\overline I$ 的复合定义算子 $G:L^2(\Omega)\to L^2(\Omega)$:
\begin{equation}\label{eq:elliptic-l2-compact-solution-operator-definition}
G:=\overline I\,\overline G,
\end{equation}
其中 $\overline I$ 为从 $H_0^k(\Omega)$ 到 $L^2(\Omega)$ 的恒等映射. 由 Theorem~\ref{thm:compact-successive-sobolev-orders} 知该算子紧. 由于有界算子与紧算子的复合为紧算子 (参见 Problem~\ref{ex:bounded-compact-operator-products}), 得到以下结论.

\begin{Lemma}
\label{lem:elliptic-l2-solution-operator-compact}
解算子 $G:L^2(\Omega)\to L^2(\Omega)$ 是紧的.
\end{Lemma}
